\documentclass[11pt]{article}

\usepackage[margin=1in]{geometry}
\usepackage{amsmath,amssymb,amsthm,mathtools}
\usepackage{mathrsfs}
\usepackage{microtype}
\usepackage{enumitem}
\usepackage{graphicx}
\usepackage[hidelinks]{hyperref}

\newtheorem{theorem}{Theorem}[section]
\newtheorem{proposition}[theorem]{Proposition}
\newtheorem{lemma}[theorem]{Lemma}
\newtheorem{corollary}[theorem]{Corollary}
\newtheorem{remark}[theorem]{Remark}
\theoremstyle{definition}
\newtheorem{definition}[theorem]{Definition}

\newcommand{\NA}{\operatorname{NA}}
\newcommand{\calL}{\mathcal L}
\newcommand{\ran}{\operatorname{ran}}
\newcommand{\rank}{\operatorname{rank}}
\newcommand{\codim}{\operatorname{codim}}
\newcommand{\R}{\mathbb R}
\newcommand{\X}{\mathfrak X}
\newcommand{\Xzero}{X_0}
\newcommand{\Yop}{\mathscr Y}
\newcommand{\norm}[1]{\left\lVert #1\right\rVert}
\newcommand{\abs}[1]{\left\lvert #1\right\rvert}
\newcommand{\sphere}[1]{S_{#1}}
\newcommand{\ball}[1]{B_{#1}}

\title{Norm-Attaining Operators of Every Finite Rank\\on Mart\'in's Space}
\author{A verification and solution report}
\date{June 22, 2026}

\begin{document}
\maketitle

\section*{Packet status}

This packet upgrades the previous partial packet for arXiv:2406.07273.  It is
classified as:
\[
  \texttt{full\_solution\_likely\_valid}.
\]
The upgraded claim is a positive solution of Mart\'in's Open Problem 7: the
space \(\X\) admits rank-two norm-attaining operators into \(\ell_2^2\).  The
argument proves the stronger finite-rank statement in Theorem \ref{thm:main}.

The earlier active packet proved only a fixed-kernel nowhere-density result and
explicitly did not settle the existence question.  That packet is preserved in
the packet history directory.  The submitted TeX and PDF report are preserved
in the packet evidence directory; the source paper and open-problem crop are
retained as source evidence.

\begin{abstract}
Miguel Mart\'in constructed a smooth renorming \(\X\) of \(c_0\) whose set of
norm-attaining functionals contains no non-trivial cone, and asked whether
\(\NA(\X,\ell_2^2)\) nevertheless contains an operator of rank two.  We first
audit a proposed fixed-kernel nowhere-density argument.  Its conclusion is
correct, but the written proof uses the false general assertion that smoothness
of a Banach space is equivalent to strict convexity of its dual.  A local
quotient-smoothness argument repairs the proof.

We then give a positive answer to Mart\'in's question.  In fact, a stronger
statement follows directly from the architecture of his renorming: for every
integer \(n\geq 1\) and every real Banach space \(Y\) with
\(\dim Y\geq n\), there is a norm-attaining operator from \(\X\) to \(Y\) of
rank \(n\).  The key observation is a general lifting lemma for a graph-sum
norm
\[
       p(x)=\norm{x}_0+\lambda\norm{Rx}.
\]
A finite-rank norm-attaining operator for the old norm can be modified by one
rank-one term in the direction of its norming value.  The triangle inequality
then gives the upper bound for the new operator norm, while equality holds at
the same vector.  Mart\'in's injectivity and range-separation properties ensure
that the rank is preserved.
\end{abstract}

\section{The question and the literature check}

Throughout, all Banach spaces and operators are real.  For Banach spaces
\(X,Y\), \(\NA(X,Y)\) denotes the set of norm-attaining bounded operators
from \(X\) to \(Y\).

Mart\'in's paper constructs a particular renorming \(\X\) of \(c_0\) and ends
with the following question:
\begin{quote}
\emph{Does \(\NA(\X,\ell_2^2)\) contain rank-two operators?}
\end{quote}
The source evidence crop from Open Problem 7 is reproduced below.

\begin{center}
\includegraphics[width=\textwidth,height=0.55\textheight,keepaspectratio]{figures/open_problem_crop.png}
\end{center}

This is Open Problem 7 in \cite{Martin2025}.  A targeted search through
June 22, 2026 found no published proof or counterexample.  In particular,
Mart\'in's May 2025 BIRS handout still states the question as open
\cite{MartinBIRS}.  Thus the argument below is presented as a new solution and
has not been peer reviewed.

Our main result is stronger than the requested rank-two assertion.

\begin{theorem}[Main theorem]\label{thm:main}
Let \(\X\) be the Banach space produced in the proof of
\cite[Theorem 1]{Martin2025}.  Let \(n\in\mathbb N\), and let \(Y\) be a real
Banach space with \(\dim Y\geq n\).  Then there exists
\(T\in\NA(\X,Y)\) with
\[
             \rank T=n.
\]
In particular, \(\NA(\X,\ell_2^2)\) contains a rank-two operator.
\end{theorem}

\section{The structural ingredients in Mart\'in's construction}

We isolate only the properties actually used below.  In Steps 1--7 of the
proof of \cite[Theorem 1]{Martin2025}, the final space \(\X\) is built as
follows.

The underlying vector space is \(c_0\).  First it is given an equivalent
smooth norm \(\norm{\cdot}_0\); write
\[
             \Xzero=(c_0,\norm{\cdot}_0).
\]
The construction has
\begin{equation}\label{eq:NA-base}
       \NA(\Xzero,\R)=\NA(c_0,\R)=c_{00}.
\end{equation}
Here the dual spaces are identified algebraically with \(\ell_1\).  Recall
that an element of \(\ell_1=c_0^*\) attains its usual norm on \(c_0\) exactly
when it has finite support, which proves the second equality in
\eqref{eq:NA-base}.

There are also a Banach space \(W\), a bounded operator
\(R:\Xzero\to W\), and a linear subspace \(\Yop\subset \Xzero^*\) such that
\begin{align}
       &R \text{ is injective},                                      \label{eq:Rinj}\\
       &R^*:W^*\longrightarrow \Xzero^* \text{ is injective},         \label{eq:Rstarinj}\\
       &R^*(W^*)\subseteq \Yop,                                      \label{eq:Rrange}\\
       &\Yop\cap\NA(\Xzero,\R)=\{0\}.                               \label{eq:separation}
\end{align}
Specifically, \eqref{eq:Rrange} and \eqref{eq:Rstarinj} are equations (3)
and (4) in Step 6 of Mart\'in's proof.  The final norm is
\begin{equation}\label{eq:final-norm}
              p(x)=\norm{x}_0+\norm{Rx}_W,
              \qquad x\in c_0,
\end{equation}
and \(\X=(c_0,p)\).  The scaled variants in
\cite[Remark 2]{Martin2025} have
\(p_\lambda(x)=\norm{x}_0+\lambda\norm{Rx}\), where \(\lambda>0\); the same
argument works for all of them.

\section{Two elementary lemmas}

We first record the finite-rank norm-attainment criterion used in
\cite[Corollary 2.6]{KLMW}.  A short proof is included.

\begin{lemma}[Finite-rank criterion]\label{lem:finite-rank-criterion}
Let \(X,Y\) be Banach spaces and let \(S:X\to Y\) have finite rank.  If
\[
             [\ker S]^\perp=\ran S^*\subseteq \NA(X,\R),
\]
then \(S\in\NA(X,Y)\).
\end{lemma}

\begin{proof}
Put \(E=\ran S\), and regard \(S\) as a surjection from \(X\) onto the
finite-dimensional space \(E\).  The unit sphere of \(E^*\) is compact, so
there is \(e^*\in S_{E^*}\) such that
\[
             \norm{S^*e^*}=\norm{S^*}=\norm{S}.
\]
The functional \(S^*e^*\) belongs to \([\ker S]^\perp\), hence attains its
norm at some \(x\in S_X\).  After changing the sign of \(x\) if needed,
\[
  \norm{S}
  =(S^*e^*)(x)
  =e^*(Sx)
  \leq \norm{Sx}
  \leq \norm{S}.
\]
Thus \(\norm{Sx}=\norm{S}\).
\end{proof}

The next lemma is the mechanism that passes from the old norm to the graph-sum
norm.

\begin{lemma}[Graph-sum lifting]\label{lem:graph-sum}
Let \((X,\norm{\cdot}_0)\) and \(W\) be Banach spaces, let
\(R:X\to W\) be bounded, and fix \(\lambda>0\).  Define
\[
          p_\lambda(z)=\norm{z}_0+\lambda\norm{Rz}.
\]
Suppose that \(S:X\to Y\) has finite rank, \(\norm{S}=1\) for the old norm,
and attains its norm at \(x\in S_{(X,\norm{\cdot}_0)}\).  Put \(y=Sx\), so
\(\norm{y}=1\).  Assume that \(Rx\neq0\), and choose
\(w^*\in S_{W^*}\) with
\[
                 w^*(Rx)=\norm{Rx}.
\]
If
\begin{equation}\label{eq:outside-range}
                 R^*w^*\notin\ran S^*,
\end{equation}
then the operator
\begin{equation}\label{eq:T-general}
        Tz=Sz+\lambda\,(R^*w^*)(z)y,
        \qquad z\in X,
\end{equation}
viewed from \((X,p_\lambda)\) to \(Y\), satisfies
\[
          T\in\NA((X,p_\lambda),Y),
          \qquad \norm{T}=1,
          \qquad \rank T=\rank S.
\]
It attains its norm at \(x/p_\lambda(x)\).
\end{lemma}

\begin{proof}
For every \(z\in X\), the triangle inequality gives
\begin{align*}
 \norm{Tz}
 &\leq \norm{Sz}+\lambda\abs{w^*(Rz)}\norm{y}\\
 &\leq \norm{z}_0+\lambda\norm{Rz}
  =p_\lambda(z).
\end{align*}
Hence \(\norm{T}\leq1\).  At the norming vector \(x\),
\[
     Tx=y+\lambda\norm{Rx}\,y
       =\bigl(1+\lambda\norm{Rx}\bigr)y,
\]
whereas
\[
     p_\lambda(x)=1+\lambda\norm{Rx}.
\]
Consequently \(\norm{T}=1\), and \(T\) attains its norm at
\(x/p_\lambda(x)\).

It remains to check the rank.  Let \(E=\ran S\); then \(y\in E\) and
\(\ran T\subseteq E\).  Regard \(S,T\) as operators into \(E\).  Since
\(S:X\to E\) is onto, \(S^*:E^*\to X^*\) is injective and
\(\ran S^*= [\ker S]^\perp\).  If \(T^*e^*=0\) for some \(e^*\in E^*\),
then
\[
       S^*e^*=-\lambda e^*(y)R^*w^*.
\]
If \(e^*(y)=0\), injectivity of \(S^*\) gives \(e^*=0\).  If
\(e^*(y)\neq0\), the displayed identity contradicts
\eqref{eq:outside-range}.  Thus \(T^*:E^*\to X^*\) is injective, so
\(T:X\to E\) is onto.  Therefore \(\rank T=\dim E=\rank S\).
\end{proof}

\section{Proof of the main theorem}

\begin{proof}[Proof of Theorem \ref{thm:main}]
Fix \(n\in\mathbb N\) and a Banach space \(Y\) with \(\dim Y\geq n\).
Choose an \(n\)-dimensional subspace \(E\subseteq Y\) and an
\(n\)-dimensional subspace
\[
                  U\subset c_{00}=\NA(\Xzero,\R).
\]
Let \(M=U_\perp=\bigcap_{f\in U}\ker f\).  Then \(\Xzero/M\) is
\(n\)-dimensional and
\[
                  M^\perp=U.
\]
Choose a linear isomorphism \(A:\Xzero/M\to E\), let
\(q:\Xzero\to\Xzero/M\) be the quotient map, and define
\[
                  S_0=\iota_E A q:\Xzero\longrightarrow Y,
\]
where \(\iota_E:E\hookrightarrow Y\) is the inclusion.  Then
\(\rank S_0=n\), \(\ker S_0=M\), and
\[
          \ran S_0^*=M^\perp=U\subseteq\NA(\Xzero,\R).
\]
By Lemma \ref{lem:finite-rank-criterion}, \(S_0\) attains its norm.  Replace
it by \(S=S_0/\norm{S_0}\), choose
\(x\in S_{\Xzero}\) with \(\norm{Sx}=1\), and put \(y=Sx\).

Since \(R\) is injective and \(x\neq0\), one has \(Rx\neq0\).  Choose
\(w^*\in S_{W^*}\) satisfying \(w^*(Rx)=\norm{Rx}\), and set
\[
                    r=R^*w^*.
\]
By \eqref{eq:Rrange}, \(r\in\Yop\).  By injectivity of \(R^*\),
\(r\neq0\).  If \(r\in\ran S^*=U\), then
\[
     r\in \Yop\cap U
       \subseteq \Yop\cap\NA(\Xzero,\R)
       =\{0\},
\]
a contradiction.  Hence \(r\notin\ran S^*\).

Lemma \ref{lem:graph-sum}, with \(\lambda=1\), now shows that
\[
          Tz=Sz+r(z)y.
\]
more simply written as
\[
          T=S+y\otimes r,
          \qquad (y\otimes r)(z)=r(z)y,
\]
is a norm-one norm-attaining operator from \((c_0,p)=\X\) to \(Y\), and
\(\rank T=\rank S=n\).  This proves the theorem.
\end{proof}

\begin{remark}[The rank-two formula]\label{rem:rank-two-formula}
For \(n=2\) and \(Y=\ell_2^2\), rotate the norm-one base operator so that
\[
                S=(f_0,h_0),\qquad Sx=(1,0).
\]
Then the resulting operator has the especially transparent form
\[
        Tz=\bigl(f_0(z)+(R^*w^*)(z),\ h_0(z)\bigr).
\]
Indeed,
\[
 \norm{Tz}_2
 \leq \norm{(f_0(z),h_0(z))}_2+\abs{w^*(Rz)}
 \leq \norm{z}_0+\norm{Rz}=p(z),
\]
with equality at \(x/p(x)\).  The two coordinate functionals are independent
because \(R^*w^*\notin\operatorname{span}\{f_0,h_0\}\).
\end{remark}

\begin{corollary}\label{cor:scaled}
The conclusion of Theorem \ref{thm:main} also holds for every scaled norm
\[
              p_\lambda(x)=\norm{x}_0+\lambda\norm{Rx},
              \qquad \lambda>0,
\]
including the arbitrarily small equivalent perturbations described in
\cite[Remark 2]{Martin2025}.
\end{corollary}

\begin{proof}
Use Lemma \ref{lem:graph-sum} with the same \(S,x,y,w^*\) and the specified
\(\lambda\).
\end{proof}

\section{Audit and repair of the fixed-kernel partial result}

The proposed partial result asserts that for every fixed two-codimensional
kernel, norm attainment is nowhere dense in the corresponding four-dimensional
operator fibre.  The conclusion is correct, but two points in the supplied
write-up need correction.

First, the finiteness statement for
\(q(B_\X)\cap S_{\X/M}\) is supplied by Proposition 6(b) of
\cite{Martin2025}, not by ``Proposition 2.4'' of that paper.  Second, the
sentence
\[
       \text{``\(\X\) is smooth; equivalently, \(\X^*\) is strictly convex''}
\]
is false for general Banach spaces.  In infinite dimensions, smoothness of a
space does not in general imply strict convexity of its dual.  Consequently,
the global claim that every two-dimensional quotient \(\X/M\) is smooth does
not follow from the stated reasoning.

Only smoothness at a finite set of quotient points is needed, and that property
follows directly from smoothness of \(\X\).

\begin{lemma}[Local smoothness passes to a represented quotient point]
\label{lem:quotient-smooth-point}
Let \(q:X\to E=X/M\) be a quotient map.  Suppose that \(x\in S_X\), that
\(e=qx\in S_E\), and that \(X\) is smooth at \(x\).  Then \(E\) is smooth at
\(e\).
\end{lemma}

\begin{proof}
Let \(\phi_1,\phi_2\in S_{E^*}\) satisfy
\(\phi_1(e)=\phi_2(e)=1\).  The adjoint \(q^*:E^*\to X^*\) is an isometry, and
\[
        (q^*\phi_i)(x)=\phi_i(qx)=1\qquad(i=1,2).
\]
Thus both \(q^*\phi_1\) and \(q^*\phi_2\) are norming functionals for \(x\).
Smoothness at \(x\) gives \(q^*\phi_1=q^*\phi_2\).  Since \(q^*\) is
injective, \(\phi_1=\phi_2\).
\end{proof}

We now give a corrected proof of the partial theorem.  The finite-dimensional
perturbation lemma needs smoothness only at the prescribed vector.

\begin{lemma}\label{lem:Fe}
Let \(E\) be a real two-dimensional Banach space, let \(H=\ell_2^2\), and
fix \(e\in S_E\).  Assume that \(E\) is smooth at \(e\).  Define
\[
        F_e=\{A\in\calL(E,H):\norm{Ae}=\norm{A}\}.
\]
Then \(F_e\) is closed and nowhere dense in \(\calL(E,H)\).
\end{lemma}

\begin{proof}
Closedness is immediate.  We show that \(F_e\) has empty interior.

Let \(A\in F_e\) and \(\delta>0\).  If \(A=0\), choose a nonzero
\(B:E\to H\) with \(Be=0\) and \(\norm{B}<\delta\); then \(B\notin F_e\).

Assume \(A\neq0\), put \(y=Ae\), and let \(j\in S_{E^*}\) be the unique
norming functional for \(e\).  Choose \(d\in\ker j\setminus\{0\}\).  The
norm of \(E\) is G\^ateaux differentiable at \(e\), so
\[
       \frac{d}{dt}\bigg|_{t=0}\norm{e+td}=j(d)=0.
\]
Set \(u=y/\norm{y}\).  Since \(\norm{A(e+td)}\leq
\norm{y}\norm{e+td}\) with equality at \(t=0\), differentiation from both
sides gives
\[
                     \langle u,Ad\rangle=0.
\]
Choose \(B\in\calL(E,H)\) with \(Be=0\) and \(Bd=y\), and put
\(A_\alpha=A+\alpha B\).  Then \(A_\alpha e=y\), while
\[
       A_\alpha(e+td)=y+tAd+\alpha t y.
\]
The derivative at zero of the numerator in
\[
       \frac{\norm{A_\alpha(e+td)}}{\norm{e+td}}
\]
is
\[
       \langle u,Ad+\alpha y\rangle=\alpha\norm{y}>0,
\]
and the derivative of the denominator is zero.  Hence, for all sufficiently
small \(t>0\),
\[
       \frac{\norm{A_\alpha(e+td)}}{\norm{e+td}}
       >\norm{y}=\norm{A_\alpha e}.
\]
Thus \(A_\alpha\notin F_e\).  Taking \(\alpha>0\) so small that
\(\norm{\alpha B}<\delta\) proves that \(F_e\) has empty interior.
\end{proof}

\begin{theorem}[Corrected fixed-kernel theorem]\label{thm:fixed-kernel}
Let \(M\) be a closed subspace of \(\X\) with \(\codim M=2\), let
\(H=\ell_2^2\), and set
\[
       \calL_M=\{T\in\calL(\X,H):M\subseteq\ker T\}.
\]
Then \(\NA(\X,H)\cap\calL_M\) is nowhere dense in \(\calL_M\).
\end{theorem}

\begin{proof}
Let \(q:\X\to E=\X/M\) be the quotient map.  The map
\(A\mapsto Aq\) identifies \(\calL(E,H)\) isometrically with \(\calL_M\).
Set
\[
                 C_M=q(B_\X)\cap S_E.
\]
By \cite[Proposition 6(b)]{Martin2025}, \(C_M\) contains at most two
antipodal pairs.

For each \(e\in C_M\), choose \(x\in B_\X\) with \(qx=e\).  Since
\(\norm{e}=1\), necessarily \(\norm{x}=1\).  The space \(\X\) is smooth, so
Lemma \ref{lem:quotient-smooth-point} shows that \(E\) is smooth at \(e\).
Lemma \ref{lem:Fe} therefore applies to every \(e\in C_M\).

We claim that
\begin{equation}\label{eq:fibre-inclusion}
 \NA(\X,H)\cap\calL_M
 \subseteq
 \{0\}\cup\bigcup_{e\in C_M}F_e.
\end{equation}
Indeed, let \(T=Aq\neq0\) attain its norm at \(x\in S_\X\).  Since
\(\norm{Aq}=\norm{A}\),
\[
                \norm{A(qx)}=\norm{A}.
\]
Writing \(e=qx\), one cannot have \(\norm{e}<1\), because then
\[
       \norm{A(e/\norm{e})}
       =\frac{\norm{Ae}}{\norm{e}}
       >\norm{A}.
\]
Thus \(e\in q(B_\X)\cap S_E=C_M\), and \(A\in F_e\).  This proves
\eqref{eq:fibre-inclusion}.  The right-hand side is a finite union of closed
nowhere dense subsets of the four-dimensional space \(\calL(E,H)\), so the
left-hand side is nowhere dense.
\end{proof}

\begin{corollary}\label{cor:fixed-kernel}
If \(T\in\NA(\X,\ell_2^2)\) has rank two and kernel \(M\), then for every
\(\varepsilon>0\) there is a rank-two operator
\(S:\X\to\ell_2^2\) such that
\[
       \ker S=M,\qquad \norm{S-T}<\varepsilon,
       \qquad S\notin\NA(\X,\ell_2^2).
\]
\end{corollary}

\begin{proof}
Write \(T=Aq\), where \(q:\X\to\X/M\) is the quotient map.  Since \(T\)
has rank two, \(A:\X/M\to\ell_2^2\) is invertible.  Invertible operators form
an open subset of \(\calL(\X/M,\ell_2^2)\).  Theorem
\ref{thm:fixed-kernel} therefore supplies, arbitrarily close to \(A\), an
invertible \(B\) such that \(Bq\) is not norm attaining.  Put \(S=Bq\).
Then \(\ker S=M\), \(\rank S=2\), and
\(\norm{S-T}=\norm{B-A}\).
\end{proof}

\section{How the two conclusions fit together}

Theorem \ref{thm:main} and Theorem \ref{thm:fixed-kernel} are fully
compatible.  The main theorem proves existence of many finite-rank
norm-attaining operators, including rank-two ones.  The fixed-kernel theorem
says that after one fixes the kernel of any such operator, norm attainment is
still a nowhere dense condition inside that particular four-dimensional
fibre.  Non-emptiness and nowhere density are not contradictory.

The argument settles only the existence question in Open Problem 7.  It does
not by itself settle the older density problem asking whether
\(\NA(X,\ell_2^2)\) is dense in \(\calL(X,\ell_2^2)\) for every Banach space
\(X\), nor whether every finite-rank operator can be approximated by
norm-attaining operators.

\section*{Conclusion}

The candidate fixed-kernel result survives audit after replacing an incorrect
duality assertion by a local quotient-smoothness lemma.  More importantly,
the graph-sum structure of Mart\'in's norm yields a direct positive solution of
Open Problem 7, and in fact produces norm-attaining operators of every
prescribed finite rank into every sufficiently large range space.

\begin{thebibliography}{9}

\bibitem{Martin2025}
M. Mart\'in,
\emph{A Banach space whose set of norm-attaining functionals is algebraically
trivial},
J. Funct. Anal. \textbf{288} (2025), 110815.
Preprint: \href{https://arxiv.org/abs/2406.07273}{arXiv:2406.07273}.

\bibitem{KLMW}
V. Kadets, G. L\'opez, M. Mart\'in, and D. Werner,
\emph{Norm attaining operators of finite rank},
in \emph{The Mathematical Legacy of Victor Lomonosov},
Advances in Analysis and Geometry 2, De Gruyter, 2020, pp. 157--188.
Preprint: \href{https://arxiv.org/abs/1905.08272}{arXiv:1905.08272}.

\bibitem{MartinBIRS}
M. Mart\'in,
\emph{A Banach space whose set of norm-attaining functionals is algebraically
trivial}, BIRS--IMAG workshop handout, May 2025.
\href{https://www.ugr.es/~mmartins/Curriculum/Charlas/2025_BIRS_handout.pdf}
{Author-hosted PDF}.

\end{thebibliography}

\end{document}
